\section{Anwendungsfälle}
\label{sec:anwendungsfaelle}

Die folgenden Fälle zeigen jeweils einen fachlich interpretierbaren Befund und
den dafür benötigten Interaktionsweg. Die drei Zustände sind keine zwingende
lineare Klickfolge: Für jeden Fall wird die passende Zeitspanne gewählt. So
führt der Dänemark-Zustand in \cref{sec:case3} nicht versehentlich den zuvor
gewählten Frankreich-Zustand fort, sondern demonstriert einen unabhängigen
Einstieg über die Matrix.

\subsection{Anwendung Visualisierung Eins}
\label{sec:case1}

Ausgangspunkt ist der 1. Mai 2025 um 00 Uhr UTC. Ein Klick auf Frankreich im
gerichteten Netzwerk wählt das Länderpaar. In \cref{fig:case-france} sind
Frankreichs Knoten gelb und seine Kante kräftig, während die übrigen Kanten
abgeblendet sind. Die rote Pfeilrichtung zeigt einen kleinen physischen Import
von Frankreich nach Deutschland (\(+0{,}058\gw\)). Der Tooltip nennt zugleich
einen Handelswert von \(-2{,}940\gw\), also kommerziellen Export in der hier
verwendeten Vorzeichenkonvention.

\begin{figure}[htbp]
  \centering
  \begin{tikzpicture}
    \node[anchor=south west,inner sep=0] (img) at (0,0)
      {\includegraphics[width=0.70\textwidth]{02_frankreich_netzwerk.png}};
    \begin{scope}[x={(img.south east)},y={(img.north west)}]
      \draw[importred,very thick,-{Latex[length=3mm]}]
        (0.91,0.68) node[above left,fill=white,rounded corners,inner sep=2pt,
          text=importred,font=\sffamily\bfseries\small]{1: gewählte Beziehung}
        -- (0.735,0.385);
      \draw[projectyellow!70!black,very thick]
        (0.735,0.38) circle[radius=0.032];
    \end{scope}
  \end{tikzpicture}
  \caption{Frankreich ist ausgewählt. Knoten und Kante bleiben hervorgehoben;
    die rote Richtung steht für physischen Import nach Deutschland.}
  \label{fig:case-france}
\end{figure}

Der Befund ist für die Zielgruppe wichtig, weil er unmittelbar zeigt, warum
Handel und physischer Fluss nicht synonym verwendet werden dürfen. Eine Tabelle
könnte die beiden exakten Werte sogar schneller liefern, würde aber nicht
gleichzeitig sichtbar machen, wie schwach die ausgewählte physische Kante im
Vergleich zu starken Exportkanten nach Österreich, in die Schweiz oder nach
Polen ist. Der Netzwerkgraph liefert den relationalen Kontext; der Tooltip
liefert die Präzision.

\subsection{Anwendung Visualisierung Zwei}
\label{sec:case2}

Für die Woche vom 8. bis 14. Mai wird Frankreich gewählt und anschlie{\ss}end der
11. Mai 2025 um 11 Uhr UTC angeklickt. Die gestrichelte Vertikale in
\cref{fig:case-price} verbindet Erzeugungsmix und Frankreich-Fluss. In dieser
Stunde beträgt die erneuerbare Leistung \(58{,}961\gw\), entsprechend
88,7\,\% der dargestellten Gesamterzeugung von etwa \(66{,}5\gw\). Der
Day-Ahead-Preis erreicht mit \(-250{,}32\mwh\) das Monatsminimum, während aus
Frankreich physisch \(+1{,}382\gw\) importiert werden.

\begin{figure}[htbp]
  \centering
  \begin{tikzpicture}
    \node[anchor=south west,inner sep=0] (img) at (0,0)
      {\includegraphics[width=0.98\textwidth]{03_negativpreis_zeitreihe.png}};
    \begin{scope}[x={(img.south east)},y={(img.north west)}]
      \draw[importred,very thick,-{Latex[length=3mm]}]
        (0.69,0.96) node[below,fill=white,rounded corners,inner sep=2pt,
          text=importred,font=\sffamily\bfseries\small]{1: gewählte Stunde}
        -- (0.535,0.69);
      \draw[projectblue,very thick,-{Latex[length=3mm]}]
        (0.87,0.90) node[below,fill=white,rounded corners,inner sep=2pt,
          text=projectblue,font=\sffamily\bfseries\small]{2: Preis und Absolutwerte}
        -- (0.535,0.17);
    \end{scope}
  \end{tikzpicture}
  \caption{Monatliches Preisminimum. Die gemeinsame Zeitmarke verbindet
    erneuerbare Erzeugung, Frankreich-Fluss und Detailwerte.}
  \label{fig:case-price}
\end{figure}

\FloatBarrier

Über den ganzen Mai korreliert der Preis mit der absoluten erneuerbaren
Leistung mit \(r=-0{,}758\), mit ihrem Anteil sogar mit \(r=-0{,}817\). Dies
ist ein sichtbares und statistisch gestütztes Monatsmuster, aber kein
Kausalnachweis. Der einfache summierte physische Fluss korreliert dagegen nur
mit \(r=0{,}084\) mit dem Preis. Die Anwendung verhindert damit eine zu einfache
Erzählung ``mehr Export erklärt negative Preise'': Der konkrete Länderfluss,
andere Erzeugungsarten, Nachfrage, Netzrestriktionen und Markteffekte können
abweichen.

Ein Scatterplot von Preis und erneuerbarer Leistung wäre für die Korrelation
präziser und ist eine sinnvolle spätere Ergänzung. Er würde jedoch die
Reihenfolge der Stunden, Tagesmuster und den zeitgleichen Frankreich-Fluss
verbergen. Für die Untersuchung eines Ereignisverlaufs ist die gekoppelte
Zeitreihe geeigneter; für einen globalen Zusammenhang wäre der Scatterplot
effizienter.

\subsection{Anwendung Visualisierung Drei}
\label{sec:case3}

Nach Rückkehr zum 48-Stunden-Fenster 3.--4. Mai wird die Matrixzelle Dänemark,
3. Mai um 12 Uhr angeklickt. \Cref{fig:case-matrix} markiert die vollständige
Dänemark-Zeile und die 12-Uhr-Spalte als Fadenkreuz. Die aktive Zelle zeigt
einen Export aus Deutschland von \(-1{,}710\gw\). In derselben Spalte sind
auch Österreich, Belgien, Tschechien, Luxemburg, Niederlande, Norwegen, Polen,
Schweden und Schweiz blau oder nahe neutral, während Frankreich rötlich ist
(\(+0{,}582\gw\) Import).

\begin{figure}[htbp]
  \centering
  \begin{tikzpicture}
    \node[anchor=south west,inner sep=0] (img) at (0,0)
      {\includegraphics[width=\textwidth]{04_daenemark_matrix.png}};
    \begin{scope}[x={(img.south east)},y={(img.north west)}]
      \draw[importred,very thick,-{Latex[length=3mm]}]
        (0.14,0.93) node[right,fill=white,rounded corners,inner sep=2pt,
          text=importred,font=\sffamily\bfseries\small]{1: Partnerzeile}
        -- (0.30,0.62);
      \draw[projectblue,very thick,-{Latex[length=3mm]}]
        (0.50,0.94) node[right,fill=white,rounded corners,inner sep=2pt,
          text=projectblue,font=\sffamily\bfseries\small]{2: Stunde}
        -- (0.346,0.62);
    \end{scope}
  \end{tikzpicture}
  \caption{Auswahl der Zelle Dänemark, 03.05.2025 12 Uhr UTC. Das goldene
    Fadenkreuz macht beide gewählten Dimensionen sichtbar.}
  \label{fig:case-matrix}
\end{figure}

\FloatBarrier

Die visuelle Aufgabe besteht nicht nur darin, einen Einzelwert abzulesen,
sondern das gleichzeitige Länderprofil zu sehen. Das Muster legt nahe, dass der
deutsche physische Austausch in dieser Stunde gegenüber vielen Partnern in
Exportrichtung verlief, aber nicht gegenüber Frankreich. Die globale Farbskala
macht die Sättigungen zeilenübergreifend vergleichbar. Eine Tabelle wäre für
den Wert \(-1{,}710\gw\) genauer, verlangte für das Profil jedoch serielles
Lesen von elf Zeilen. Elf Linien hätten stärkere Überdeckung. Die Matrix
unterstützt das parallele Erkennen, die gekoppelten Detailansichten übernehmen
die quantitative Prüfung.

\FloatBarrier
